\section{Introduction}

Numerical weather prediction (NWP) launches a forecast by integrating a numerical
model of the atmosphere forward from an estimate of its current state. Because the
atmosphere is observed only incompletely, indirectly, and noisily, producing that
initial estimate is itself an inverse problem: given a forecast model and a stream of
imperfect measurements distributed in space and time, recover the system state well
enough to forecast from it. This is the task of data assimilation (DA), and it is the
component of an operational NWP system that most directly limits forecast skill
\citep{kalnay2003atmospheric}. The flagship variational method is
four-dimensional variational data assimilation (4D-Var)
\citep{ledimet1986variational, courtier1994strategy}, which estimates the initial
condition $\xstate_0$ of an assimilation window by minimizing a cost that balances a
prior (``background'') estimate against the misfit between the model trajectory and the
observations spread across that window.

DA is hard precisely because the systems of interest are chaotic: small errors in the
analysis (the assimilated initial state) grow exponentially, so the analysis must be
accurate and the optimization landscape is treacherous. The Lorenz-96 model
\citep{lorenz1996predictability} is the canonical low-order testbed for exactly this
difficulty, and it is the primary one-dimensional (1D) system we study here. In our
configuration (grid size 40, forcing $F=8$), trajectories started $10^{-6}$ apart
separate at a fitted exponential rate $\lambda = 2.263\,/\mathrm{MTU}$, corresponding to
a Lyapunov time $T_L = 1/\lambda = 0.442$ model time unit (MTU); this is somewhat
shorter than the $\sim$0.6 MTU expected from the original study, a known artifact of our
fixed-step integrator that we report honestly. Because errors saturate within roughly a
Lyapunov time, an inaccurate analysis cannot be rescued downstream, and this is the
regime in which a better treatment of the observations is supposed to help.

\paragraph{The learned inverse-observation idea.}
\citet{frerix2021invobs} study the variational problem of recovering an initial state
$\xstate_0$ whose model trajectory $\Mop^t(\xstate_0)$ fits a sequence of observations.
The observation operator $\Hop$ in their setting (and ours) is a sparse spatial
subsampling: it keeps every fourth grid point for Lorenz-96 (10 of 40 points observed)
and every sixteenth point in each direction for two-dimensional (2D) Kolmogorov flow (a
$4\times4=16$ of 4096 points observed). Because $\Hop$ is many-to-one and highly lossy,
the standard observation-space objective
$\lVert \Hop\,\Mop^t(\xstate_0) - \yobs_t\rVert^2$ is strongly non-convex and riddled
with bad local minima, so a poorly initialized gradient-based optimizer gets trapped.
Their contribution is to train a convolutional neural network (CNN) that acts as a
learned inverse observation operator $\Hinv$, mapping an observation sequence back into
physics (state) space. They use it in two complementary ways: as a better
\emph{initialization} (lift the observations to a physics-space first guess
$\Hinv(Y)$ instead of a naive interpolation or average), and as a \emph{reformulation of
the loss} (replace the observation-space misfit with a physics-space misfit
$\sum_t \lVert \Mop^t(\xstate_0) - [\Hinv(Y)]_t\rVert^2$, which removes the lossy $\Hop$
from inside the objective and is therefore much closer to convex). They further propose
an optional \emph{hybrid} schedule: run the smoother physics-space optimization first as
a warm start, then refine against the true observations in observation space. Crucially
for our work, the original study assimilates \emph{clean} observations and uses a
\emph{simplified} objective (the observation-misfit term only, with no background term
and no explicit observation-error covariance), and its reference implementation is in
JAX/Flax.

\paragraph{Our contribution.}
This capstone tests whether the learned inverse-observation idea survives a more
operationally realistic setting, through four changes. (i)~\emph{Full 4D-Var cost.}
We replace the paper's observation-misfit-only objective with the complete 4D-Var cost
$J(\xstate_0) = \Jb + \Jo$,
\[
  J(\xstate_0) = \tfrac12 (\xstate_0-\xstate_b)^\top \Bcov^{-1} (\xstate_0-\xstate_b)
  + \tfrac12 \sum_{t} \big(\Hop \Mop^t(\xstate_0) - \yobs_t\big)^\top \Rcov^{-1}
    \big(\Hop \Mop^t(\xstate_0)-\yobs_t\big),
\]
with background covariance $\Bcov = \sigb^2 \bm{C}$ ($\bm{C}$ the system's spatial
correlation) and observation-error covariance $\Rcov = \sigobs^2 \bm{I}$. The background
term pulls the analysis toward a prior $\xstate_b$ and is what makes the problem
well-posed when observations are sparse; $\sigb$ is the background-error standard
deviation we tune and $\sigobs$ the observation-error standard deviation.
(ii)~\emph{Observation noise.} Every experiment adds independent Gaussian observation
noise with $\sigobs \in \{0.1, 0.5, 1.0\}$ (the paper used clean observations), which
both stress-tests robustness and probes a deliberate mismatch, since the inverter is
trained at one noise level and evaluated at others. (iii)~\emph{Sliding-window cycling.}
We add an operational sliding-window (cycling) 4D-Var extension in which short
assimilation windows are advanced through time and each cycle's analysis is propagated
forward to form the background for the next window---the way 4D-Var is actually run
operationally, in contrast to the paper's single fixed window.
(iv)~\emph{PyTorch reimplementation.} We rebuild the entire pipeline in PyTorch
\citep{paszke2019pytorch}---the dynamical-system integrators, the CNN inverter, the
correlation preconditioning, and the data-assimilation loop solved with the
limited-memory Broyden--Fletcher--Goldfarb--Shanno (L-BFGS) optimizer
\citep{liu1989limited}.

Our headline finding is that the paper's central claim broadly holds under these harder
conditions: learned-inverse-observation initialization combined with hybrid optimization
attains the lowest analysis root-mean-square error (RMSE) at every tested noise level,
and hybrid optimization dramatically rescues a cheap climatological-mean baseline that
observation-only optimization leaves stuck in a bad local minimum. We also report honest
negative results, notably that a secondary expectation about how the optimal $\sigb$
should shift did not hold, and that on the smoke-scale Kolmogorov-flow runs the cycling
extension does not yet show a decisive advantage.

\paragraph{Roadmap.}
The remainder of the paper is organized as follows. Section~2 summarizes the learned
inverse-observation approach of \citet{frerix2021invobs} and states precisely how we
extend it. Section~3 describes the models and methods: the Lorenz-96 and Kolmogorov-flow
dynamical systems and their integrators, the observation-noise model, the full 4D-Var
cost and its physics-space hybrid companion, the learned inverse observation operator,
the decorrelation preconditioning, and the sliding-window cycling scheme. Section~4
presents results, first for Lorenz-96 (integrator diagnostics, inverter training,
baseline initialization, the tuned-$\sigb$ full 4D-Var sweep, and the sliding-window
experiments) and then for two-dimensional Kolmogorov flow. Section~5 discusses
limitations and caveats, separating what reproduces the paper from what our extensions
add, and concludes.
